\section{Introduction}
\label{sec:intro}

\IEEEPARstart{A}{} reconfigurable intelligent surface (RIS) can extend coverage to blocked areas of a building~\cite{Basar2019Access,DiRenzo2020JSAC,Wu2021TCOM}; a 256-element RIS has served as a transmit antenna~\cite{Dai2020Access}, and another prototype has strengthened a link through a concrete wall~\cite{Pei2021TCOM}. Its benefit, however, depends on configuring the passive elements for the user's current channel rather than for a past one. Choosing those phases is a hard, non-convex problem, attacked with iterative solvers~\cite{Wu2019TWC,Abeywickrama2020TCOM,Pan2020TWC,Guo2020TWC,Choi2024TCOM} and, increasingly, with learned policies~\cite{Huang2020JSAC}. These methods are usually compared by the rate they reach on the channel they were given.

Such a comparison neglects the time the optimization itself requires. A pedestrian user changes the radio channel within tens of milliseconds, whereas an iterative solver may require seconds; when its solution becomes available, the channel for which it was computed has already changed. A coarse phase configuration delivered within the coherence time can therefore outperform a near-optimal one delivered later, and additional iterations can reduce, rather than increase, the achieved rate.

This paper quantifies this effect. We evaluate every method on the channel that exists when it finishes, its \emph{effective rate}, with latencies measured on a CPU, in an indoor office with walking users and the 3GPP TR~38.901 channel model~\cite{3gpp_channel}. We address the question a base station (BS) faces in this setting: once computation time is accounted for, which method should it run, and for how many iterations? It continues our single-pass phase predictor with a sine--cosine representation~\cite{SavasciSen2026AISY} by examining when such speed is beneficial once run time is charged.

\subsection{Contributions}
\begin{enumerate}
    \item We introduce the effective rate, which evaluates each optimizer on the channel at its measured completion time (Section~\ref{sec:framework}), and show that, for pedestrian users at the measured step cost, the solvers with the highest initial-channel rate are outperformed by a single-user heuristic that requires milliseconds.
    \item We show that the effective rate has an interior maximum in the number of iterations, governed by a single quantity, the step time relative to the Doppler frequency. A small-delay model (Proposition~1) shows that an interior optimum exists for every nonzero step time and, with fitted constants, explains how it scales, so the step count should be adapted to the user speed and the step time. In the narrowband model the optimum persists for a second classical optimizer, a larger surface and every standardized indoor speed.
    \item For continuous re-optimization at the BS, we show that the pilot overhead determines the update period and that updates of at most 8 steps suffice.
    \item We compare a learned update rule with a tuned classical optimizer both at equal steps and at equal measured latency, and show that its per-step advantage is reversed once each step is charged its measured time.
    \item We quantify the dependence of the results on the assumption that the BS refreshes its precoder from the current channel, and evaluate the effects of pilot airtime, estimation error and discrete phases.
\end{enumerate}
